\documentclass[11pt]{article}
\usepackage[margin=1in]{geometry}
\usepackage{booktabs}
\usepackage{hyperref}
\usepackage{enumitem}
\usepackage{xcolor}

\title{Reversible Agency: A Permissioned, Provenance-Carrying\\
Action Layer for LLM Operating Systems}
\author{Youssef Hariri\\
DeeWaanAI\\
\texttt{github.com/DeeWaanAI}}
\date{Draft --- October 2026}

\newcommand{\sys}{Khatt}

\begin{document}
\maketitle

\begin{abstract}
LLM agents can now call tools, move files, and send messages. The literature has
made them good at \emph{acting}. It has not made them accountable for it. We
present \sys{}, the action layer of the DeeWaanAI OS stack, built on one claim:
every agent action must be a \textbf{bound record} --- what was asked, what was
done, on whose authority, with what data --- and every action must be
\textbf{reversible} or explicitly marked irreversible before it runs. In \sys{},
an action cannot execute unless it produces this record; permission scopes are
checked by deterministic code at commit time, never by model judgment; and an
executable undo recipe is captured before execution, not reconstructed after.
Over a frozen battery of 41 actions, \sys{} restored state byte-for-byte for
100\% of undoable actions, correctly gated both irreversible actions (email send,
trash purge) behind explicit recorded confirmation, caught 15/15 injected
permission and precondition defects before execution with zero side effects, and
produced records
from which every action was fully re-narratable, with tampering to the action log
detected by a hash chain. We also report what the system provably cannot do:
irreversible actions cannot be undone, only gated; undo conflicts after external
modification are detected and refused, not resolved; and the gate constrains
normalized actions, not the fidelity of intent parsing itself.
\end{abstract}

\section{Introduction}

An LLM operating system becomes useful --- and dangerous --- at the point where it
stops rendering and starts acting. Moving a file, booking a meeting, or sending
an email differs from generating a view in one decisive way: the world changes,
and the change may not come back.

Current agent systems treat an action as a transient function call. Arguments go
in, a result comes out, and benchmarks score whether the task succeeded. What is
missing is the receipt: a durable record of what was asked, who authorized it,
what data it touched, and how to take it back. Operating systems solved this
problem for programs decades ago with permissions, capabilities, and audit logs.
The emerging LLM-OS literature reuses the vocabulary --- kernels, processes,
scheduling --- while omitting the enforcement and accountability core.

\sys{} is the action layer of DeeWaanAI OS, a layered architecture in which each
layer is simultaneously a software module and an academic paper. Layer~1, Qafiya,
provides a provenance-first data substrate in which every record carries its
origin. \sys{} extends that discipline from \emph{data} to \emph{action}. Its
contribution is not a new tool-calling mechanism, a new permission filter, or a
new logging format. It is the integration of three properties into one artifact:

\begin{enumerate}[nosep]
  \item \textbf{Permissioned}: no action executes outside an explicit,
        structured grant, checked by deterministic code at commit time.
  \item \textbf{Recorded}: every action --- including refused and failed ones ---
        is stored as a provenance-carrying Qafiya record in a hash-chained log.
  \item \textbf{Reversible}: an executable undo recipe is captured before
        execution, and every action is classified \texttt{EXACT},
        \texttt{COMPENSABLE}, or \texttt{IRREVERSIBLE} before the user commits.
\end{enumerate}

The design rule that ties them together: \emph{an action that cannot produce a
receipt cannot run}.

\section{Related work}

A structured review (805 deduplicated records; 38 verified works; OpenAlex and
arXiv, searched October 2026) maps six literatures. We summarize the synthesis;
the full review with verified citations accompanies this paper.

\textbf{Tool use.} Toolformer~\cite{toolformer}, ReAct~\cite{react},
Gorilla~\cite{gorilla}, ToolLLM~\cite{toolllm}, and HuggingGPT~\cite{hugginggpt}
established invocation, with surveys consolidating the
field~\cite{toolsurvey,agenticsurvey}. A tool call in this literature is a
transient invocation: intent, authority, and data lineage are not part of the
call's record. The field optimized the verb and discarded the receipt.

\textbf{Agent operating systems and protocols.} AIOS~\cite{aios} and the
LLM-as-OS vision~\cite{llmos} apply OS vocabulary to scheduling and context;
MCP standardizes the wire format of tool calls, with its security landscape now
actively mapped~\cite{mcplandscape}. None standardizes a permission grant, an
action record, or an undo channel.

\textbf{Benchmarks.} OSWorld~\cite{osworld}, WebArena~\cite{webarena},
$\tau$-bench~\cite{taubench}, WorkArena~\cite{workarena}, and
AgentDojo~\cite{agentdojo} score task completion and attack resistance. No
surveyed environment measures recourse: whether state can be restored, or whether
a third
party can reconstruct why an action was permitted. Reversibility and
auditability are unmeasured --- and therefore unresearched --- properties.

\textbf{Permissioning and injection defense.} CaMeL~\cite{camel} defeats prompt
injection by design with capability-style policies; MCP-Secure~\cite{mcpsecure},
MiniScope~\cite{miniscope}, and Permission Manifests~\cite{permmanifests}
implement runtime access control, least privilege, and declared scopes.
PrivacyAsst~\cite{privacyasst} guards privacy in tool use. The counterweight is
sobering: adaptive attacks break deployed defenses~\cite{adaptiveattacks}. Two
consequences follow for design: permission must be a structural field checked at
commit time, not an advisory classifier; and containment must be evaluated
adversarially, expecting bypasses.

\textbf{Reversibility.} The saga pattern~\cite{saga} pairs forward actions with
compensating transactions; SagaLLM~\cite{sagallm} ports transaction guarantees to
multi-agent LLM planning; reversible-systems research formalizes
invertibility~\cite{reversiblesurvey}. Compensation is proven --- but applied to
workflow failure, not user recourse, and it assumes hand-written compensations.
No surveyed system derives and stores a per-action undo recipe at execution time,
or classifies irreversibility for the user before commitment.

\textbf{Audit and oversight.} AgentTrace~\cite{agenttrace},
PROV-AGENT~\cite{provagent}, provenance-sensitivity audits~\cite{provsensitivity},
and pre-action authorization~\cite{preactionauthz} capture fragments: what
happened, what data flowed, who consented. Human-in-the-loop
systems~\cite{magenticui,actioncard} add approval, itself an attack
surface~\cite{loopjacking}; oversight limits are argued broadly~\cite{oversight}.
None fuses the three into a single record sufficient to re-narrate the action
without rerunning anything.

\textbf{Positioning.} \sys{}'s closest neighbours are SagaLLM (transactions for
planning, no user-facing reversibility classes), CaMeL (capability defense at the
model, no durable record), and PROV-AGENT (provenance of interactions, no undo).
The gap \sys{} fills: permission, record, and undo as one artifact in the data
substrate.

\section{Design}

\subsection{The action record}

Every action, from proposal to undo, is one record in the Qafiya store. It
carries: the user's verbatim intent and its normalized form; the actor (agent and
session); the authority (user and grant); an embedded copy of the permission
scope; machine-checked preconditions with results; the reversibility class; the
undo recipe; data inputs with pre/post content hashes; the execution log; the
outcome (\texttt{success}, \texttt{failed}, \texttt{refused},
\texttt{partial}); undo status; and recorded confirmations. Refusals and failures
are terminal states with kept records: negative results are data, not noise.

\subsection{The permission gate}

Permission grants are immutable records: action types, a target pattern,
constraints, and expiry. The gate checks, in order and in deterministic code:
grant exists and unexpired; action type is granted; every target matches the
pattern \emph{after path canonicalization}, so traversal attempts such as
\texttt{../} fail by construction; constraints hold; and actions requiring
confirmation carry a recorded confirmation before execution. Any failure refuses
the action with a machine-readable reason. No model judgment is consulted: the
adaptive-attack evidence shows advisory filters degrade under pressure, so the
enforcement core is structural.

\subsection{Undo semantics}

Undo recipes are executable operations captured \emph{before} execution: source
paths and content snapshots for file edits; trash identifiers for deletes (captured
at the moment of the move, which is what creates them); counter-operations for
calendar and draft actions. Three classes are assigned at
arm time and shown to the user before commitment:

\begin{itemize}[nosep]
  \item \texttt{EXACT} --- undo restores byte-for-byte state, verified by
        content-hash comparison.
  \item \texttt{COMPENSABLE} --- a counter-action restores semantic state
        (trash restore, field restore); the undo is itself a chained record.
  \item \texttt{IRREVERSIBLE} --- no undo exists (a sent email, a purged trash
        entry). These actions require an explicit recorded confirmation and
        carry the provider's message identifier as provenance of the external
        effect.
\end{itemize}

If a recipe cannot be captured, the action is classified \texttt{IRREVERSIBLE}
automatically. Deletes never unlink: they move entries into a managed trash
inside the Qafiya store, where trash contents are themselves provenance-carrying
records, and purging the trash is a recorded irreversible action.

Undo is verified, not assumed. Conflict detection compares current content hashes
against the record's post-action hashes: if the world changed since the action,
undo refuses with a reason rather than silently overwriting a user's later work.

\subsection{Tamper-evident log}

The action log is a hash-chained history inherited from the Qafiya substrate
(v1.3)~\cite{qafiya}: every state transition of every record appends a chained
entry to the dataset's \texttt{history.jsonl}, so modification or deletion of
history breaks the chain and is detected by a verification pass that recomputes
every hash. The mechanism is proven at the substrate and applied at the action
layer. Tamper-evidence detects altered
history; it does not prevent wholesale deletion of the store, which we leave to
replication and backup as future work.

\section{Implementation}

The prototype is a Python package implementing the lifecycle state machine
(\textsc{proposed} $\to$ \textsc{gated} $\to$ \textsc{armed} $\to$
\textsc{executing} $\to$ \textsc{done}/\textsc{failed}/\textsc{refused} $\to$
\textsc{undone}), the deterministic gate, and three adapters: a file executor
over Qafiya stores, a calendar adapter, and an email adapter that hands delivery
to an existing sending system (we integrate; we do not rebuild). The record is
written at every state transition; an action whose record cannot be written does
not run. The Layer-2 view renderer is mocked behind a narrow interface
(\texttt{render\_confirmation}, \texttt{render\_audit\_view}) until its frozen
spec exists; the engine depends on the interface, not the mock. Code is AGPL-3.0
under the DeeWaanAI GitHub organization.

\section{Evaluation}

Experiments follow the Qafiya discipline: fixtures frozen before runs, and a
spine script (\texttt{run\_all.py}) that recomputes every headline number from
raw artifacts. If the script fails, the number does not appear here. Fixtures: a
frozen battery of 41 actions across fifteen action types, each executed against a
fresh, identically seeded store, and 15 injected defects. All numbers below are
script output.

\subsection{Track A: reversibility}

A framing note before the numbers. Exact reversibility of \texttt{EXACT} actions
is true \emph{by construction}: the undo recipe is captured before execution, so
restoring it should always work. We therefore read the 37/37 figure as a
correctness check on the implementation --- a floor, not a ceiling --- not as the
finding. The informative results are the three behaviors the surveyed literature
does not offer: reversibility is \emph{classified and shown to the user before
commitment}; undo is \emph{refused when the world changed} after the action; and
irreversible actions are \emph{gated behind explicit recorded confirmation}.

\begin{table}[h]
\centering
\begin{tabular}{@{}lr@{}}
\toprule
Metric & Result \\
\midrule
Battery size & 41 actions \\
Undoable actions (EXACT + COMPENSABLE) & 37 \\
Restored byte-for-byte (state-hash verified) & \textbf{37/37 = 100\%} \\
Irreversible actions correctly gating undo & 2/2 (email send, trash purge) \\
Limit probes passed & 2/2 \\
All statuses and classes as specified & yes \\
\bottomrule
\end{tabular}
\caption{Track A: reversibility over the frozen action battery. The 37/37 figure
is a by-construction correctness check; the findings are pre-commitment
classification, conflict refusal, and irreversibility gating.}
\end{table}

Undo requested twice is a
recorded no-op, not an error. Undo after an external modification is detected via
hash mismatch and \emph{refused with a reason} --- the system protects the user's
later change instead of destroying it. Irreversible actions behave honestly: the
sent email reports \texttt{undo: impossible}, and its record carries the provider
message ID as the audit trail of an effect that cannot be taken back.

\subsection{Track B: permission containment}

Fifteen defects were injected in the style of the Qafiya integrity track:
out-of-scope targets, three path-traversal variants, an expired grant, a missing
grant, an unknown action type, failed preconditions (ghost files, ghost events,
malformed recipients, occupied destinations), missing irreversible-action
confirmations, and a two-step chain in which a legal in-scope edit is followed by
a quiet out-of-scope edit.

\begin{table}[h]
\centering
\begin{tabular}{@{}lr@{}}
\toprule
Metric & Result \\
\midrule
Injected defects & 15 \\
Caught before execution & \textbf{15/15 = 100\%} \\
State change on refused actions & none \\
\bottomrule
\end{tabular}
\caption{Track B: defect-injection containment.}
\end{table}

The two-step chain deserves note: the legal first action executed and was
recorded; the out-of-scope second action was refused at the gate, and the store
state reflected only the legal action. Containment is per-action and structural,
not per-session and statistical.

\subsection{Track C: auditability}

\begin{table}[h]
\centering
\begin{tabular}{@{}lr@{}}
\toprule
Metric & Result \\
\midrule
Actions fully re-narratable from record alone & \textbf{39/39} \\
Hash chain verified over full run & yes \\
Tamper by record modification detected & yes \\
Tamper by record deletion detected & yes \\
\bottomrule
\end{tabular}
\caption{Track C: auditability and tamper-evidence.}
\end{table}

Re-narratability is the strong property: a complete account --- intent,
authority, scope, preconditions, steps, outcome, undo path --- reconstructed from
the record alone, with no reruns and no side logs. Two battery actions were
skipped in the full-chain run because earlier actions consumed their targets;
this is recorded in the results artifact, not hidden.

\subsection{A worked example}

A request to ``fix the Q3 numbers'' produces, in order: a record with the
verbatim intent and its normalization (\texttt{file.edit},
\texttt{reports/q3.txt}); a gate check against the embedded grant; a precondition
pass (\texttt{target\_exists}); classification \texttt{EXACT}; capture of a
snapshot-based undo recipe; execution; and a final record carrying pre/post
content hashes. An undo request replays the recipe and the post-state hashes
equal the pre-state hashes --- the system does not claim the undo worked; it
proves it. Contrast with the refusal path: ``wipe everything'' under a grant
limited to \texttt{file.edit} in \texttt{reports/**} never reaches the executor,
and its refusal record is itself part of the audit trail.

\section{What \sys{} provably cannot do}

Honesty is architecture, not rhetoric. The limits, measured where possible:

\begin{enumerate}[nosep]
  \item \textbf{Irreversible actions cannot be undone.} A sent email is gone;
        \sys{} gates it, records it, and says so before commitment --- nothing
        more is possible, and the system claims nothing more.
  \item \textbf{Undo conflicts are detected and refused, not resolved.} When the
        world changed after an action, automatic undo would destroy newer state.
        Resolution policy is future work.
  \item \textbf{The gate constrains normalized actions, not intent fidelity.} If
        the parser misunderstands the request, the record faithfully documents a
        correctly-gated wrong action. Intent quality is Layer 5's problem; this
        paper claims enforcement and recourse, not understanding.
  \item \textbf{Tamper-evidence is not tamper-proofing.} The chain detects
        altered or deleted history; it does not prevent destruction of the
        whole store.
  \item \textbf{Containment is evaluated against our injection set.} Fifteen
        defects, adversarially chosen, is a floor, not a ceiling; the adaptive
        attack literature warns that static defense claims age poorly, so our
        harness is built to be extended.
\end{enumerate}

\section{AI-usage disclosure}

Kimi K3 (Moonshot AI) was used as the engineering and drafting assistant for this
layer: it implemented the prototype and experiment harness under the frozen
specification, ran the verified literature search, and produced draft text that
the author reviewed, corrected, and approved. No model made unsupervised
decisions: the specification, the frozen fixtures, the three owner design
decisions (email via an external sending system; managed trash inside the Qafiya
store; hash-chained log from v1), and this paper's claims are the author's.

AI involvement is itself audited, not asserted. Every design decision is logged
(D1--D6 in the accompanying decision log), and every number in this paper traces
through a two-way claim map to a frozen artifact; the semantic map checks that
the prose means what the data says. Three oversight episodes are on the record:

\begin{enumerate}[nosep]
  \item \textbf{E1.} The first spine run reported 64.9\% reversibility. The
        defect was in the AI-written undo code (empty parent directories left
        behind) and in the AI-written measurement (byte-wise rather than
        semantic hashing of adapter state). Caught by the frozen-spine
        discipline before any number reached the paper; fixed; re-run to
        37/37. \emph{Lesson: headline numbers must be recomputed from raw
        artifacts by a script that can fail.}
  \item \textbf{E2.} The draft claimed ``nine action types''; the frozen battery
        contains fifteen --- a paper-only number introduced during AI drafting.
        Caught by the numeric claim map and corrected. \emph{Lesson: a number
        that exists only in the paper is a defect.}
  \item \textbf{E3.} The draft said ``two battery actions were skipped'' beside
        a 36/36 table; the artifact said five. Root cause was an AI-written
        harness bug marking artifact-less setup actions as failed dependencies.
        Caught by the semantic pass comparing prose against frozen data; harness
        fixed; re-run gives two legitimate skips and 39/39. \emph{Lesson: value
        fidelity is not meaning fidelity; interpretive sentences need their own
        evidence links.}
\end{enumerate}

The enforcement artifacts --- the spine script, the frozen fixtures, the
two-way verifier, and both claim maps --- ship with this paper, so a reviewer
can re-run every check rather than take this disclosure on faith.

\section{Conclusion}

The agent literature learned to act. \sys{} asks the preceding questions: may it
act, and can it take it back? By binding permission, record, and undo recipe into
one provenance-carrying artifact in the Qafiya substrate, the action layer makes
every action auditable from its own record and reversible by construction where
physics allows --- and honest where it does not. The frozen experiments show the
design holds: 100\% exact reversibility for undoable actions, 100\% containment
of injected permission and precondition defects, and complete re-narratability under a
tamper-evident chain. The remaining layers of DeeWaanAI OS --- memory, routing,
and the shell --- build on this record as their contract.

\begin{thebibliography}{99}

\bibitem{qafiya} Hariri, Y. (2026). Qafiya: A Provenance-First Data Model for
LLM-Mediated Personal Computing (DeeWaanAI OS Layer 1). DeeWaanAI. DOI
placeholder pending publication; spec v1.1 frozen, hash-chained history v1.3.

\bibitem{toolformer} Schick, T., Dwivedi-Yu, J., Dessì, R., et al. (2023).
Toolformer: Language Models Can Teach Themselves to Use Tools. arXiv.
\url{https://doi.org/10.48550/arXiv.2302.04761}

\bibitem{react} Yao, S., Zhao, J., Yu, D., et al. (2022). ReAct: Synergizing
Reasoning and Acting in Language Models. arXiv.
\url{https://doi.org/10.48550/arXiv.2210.03629}

\bibitem{gorilla} Patil, S. G., Zhang, T., Wang, X., et al. (2023). Gorilla:
Large Language Model Connected with Massive APIs. arXiv.
\url{https://doi.org/10.48550/arXiv.2305.15334}

\bibitem{toolllm} Qin, Y., Liang, S., Ye, Y., et al. (2023). ToolLLM:
Facilitating Large Language Models to Master 16000+ Real-world APIs. arXiv.
\url{https://doi.org/10.48550/arXiv.2307.16789}

\bibitem{hugginggpt} Shen, Y., Song, K., Tan, X., et al. (2023). HuggingGPT:
Solving AI Tasks with ChatGPT and its Friends in Hugging Face. arXiv.
\url{https://doi.org/10.48550/arXiv.2303.17580}

\bibitem{toolsurvey} Qu, C., Dai, S., Wei, X., et al. (2025). Tool learning with
large language models: a survey. Frontiers of Computer Science.
\url{https://doi.org/10.1007/s11704-024-40678-2}

\bibitem{agenticsurvey} Plaat, A., van Duijn, M. J., van Stein, B., et al.
(2025). Agentic Large Language Models, a Survey. JAIR.
\url{https://doi.org/10.1613/jair.1.18675}

\bibitem{aios} Mei, K., Zhu, X., Xu, W., et al. (2024). AIOS: LLM Agent
Operating System. arXiv. \url{https://doi.org/10.48550/arXiv.2403.16971}

\bibitem{llmos} Ge, Y., Ren, Y., Hua, W., et al. (2023). LLM as OS, Agents as
Apps. arXiv. \url{https://doi.org/10.48550/arXiv.2312.03815}

\bibitem{mcplandscape} Hou, X., Zhao, Y., Wang, S., et al. (2026). Model Context
Protocol (MCP): Landscape, Security Threats, and Future Research Directions.
ACM TOSEM. \url{https://doi.org/10.1145/3796519}

\bibitem{osworld} Xie, T., Zhang, D., Chen, J., et al. (2024). OSWorld:
Benchmarking Multimodal Agents for Open-Ended Tasks in Real Computer
Environments. arXiv. \url{https://doi.org/10.48550/arXiv.2404.07972}

\bibitem{webarena} Zhou, S., Xu, F. F., Zhu, H., et al. (2023). WebArena: A
Realistic Web Environment for Building Autonomous Agents. arXiv.
\url{https://doi.org/10.48550/arXiv.2307.13854}

\bibitem{taubench} Yao, S., Shinn, N., Razavi, P., et al. (2024). $\tau$-bench:
A Benchmark for Tool-Agent-User Interaction in Real-World Domains. arXiv.
\url{https://doi.org/10.48550/arXiv.2406.12045}

\bibitem{workarena} Drouin, A., Gasse, M., Caccia, M., et al. (2024). WorkArena:
How Capable Are Web Agents at Solving Common Knowledge Work Tasks? arXiv.
\url{https://doi.org/10.48550/arXiv.2403.07718}

\bibitem{agentdojo} Debenedetti, E., Zhang, J., Balunović, M., et al. (2024).
AgentDojo: A Dynamic Environment to Evaluate Prompt Injection Attacks and
Defenses for LLM Agents. arXiv. \url{https://doi.org/10.48550/arXiv.2406.13352}

\bibitem{camel} Debenedetti, E., Shumailov, I., Fan, T., et al. (2025).
Defeating Prompt Injections by Design (CaMeL). arXiv.
\url{https://doi.org/10.48550/arXiv.2503.18813}

\bibitem{mcpsecure} Singh, G., Madisetti, V. K. (2026). MCP-Secure: A Runtime
Access Control Layer for Privilege-Aware LLM Agent Tooling. IEEE OJCS.
\url{https://doi.org/10.1109/ojcs.2026.3664314}

\bibitem{miniscope} Zhu, J., Huang, X., Tseng, K., et al. (2025). MiniScope:
Authorizing Agents with Least-Privilege Permissions. arXiv.
\url{https://doi.org/10.48550/arXiv.2512.11147}

\bibitem{permmanifests} Marro, S., Chan, A., Ren, X., et al. (2026). Permission
Manifests for Web Agents. arXiv. \url{https://doi.org/10.48550/arXiv.2601.02371}

\bibitem{adaptiveattacks} Zhan, Q., Fang, R., Panchal, H. S., et al. (2025).
Adaptive Attacks Break Defenses Against Indirect Prompt Injection Attacks on
LLM Agents. NAACL Findings. \url{https://doi.org/10.18653/v1/2025.findings-naacl.395}

\bibitem{privacyasst} Zhang, X., Xu, H., Ba, Z., et al. (2024). PrivacyAsst:
Safeguarding User Privacy in Tool-Using Large Language Model Agents. IEEE TDSC.
\url{https://doi.org/10.1109/tdsc.2024.3372777}

\bibitem{saga} Štefanko, M., Chaloupka, O., Rossi, B. (2019). The Saga Pattern
in a Reactive Microservices Environment. ICSOFT.
\url{https://doi.org/10.5220/0007918704830490}

\bibitem{sagallm} Chang, E. Y., Geng, L. (2025). SagaLLM: Context Management,
Validation, and Transaction Guarantees for Multi-Agent LLM Planning. PVLDB.
\url{https://doi.org/10.14778/3750601.3750611}

\bibitem{reversiblesurvey} Mezzina, C. A., Schlatte, R., Glück, R., et al.
(2020). Software and Reversible Systems: A Survey of Recent Activities. LNCS.
\url{https://doi.org/10.1007/978-3-030-47361-7_2}

\bibitem{agenttrace} AlSayyad, A., Huang, K. Y., Pal, R. (2026). AgentTrace: A
Structured Logging Framework for Agent System Observability. arXiv.
\url{https://doi.org/10.48550/arXiv.2602.10133}

\bibitem{provagent} de Souza, R. P., Gueroudji, A., DeWitt, S. J., et al.
(2025). PROV-AGENT: Unified Provenance for Tracking AI Agent Interactions in
Agentic Workflows. IEEE eScience.
\url{https://doi.org/10.1109/escience65000.2025.00093}

\bibitem{provsensitivity} Liao, J. (2026). Auditing Provenance Sensitivity in
LLM Agent Action Selection. arXiv.
\url{https://doi.org/10.48550/arXiv.2607.20827}

\bibitem{preactionauthz} Grice, D. (2026). Provenance-Based Pre-Action
Authorization for LLM Agents. Zenodo.
\url{https://doi.org/10.5281/zenodo.21270299}

\bibitem{magenticui} Mozannar, H., Bansal, G., Tan, C., et al. (2025).
Magentic-UI: Towards Human-in-the-loop Agentic Systems. arXiv.
\url{https://doi.org/10.48550/arXiv.2507.22358}

\bibitem{actioncard} Irshad, H., Mughees, A., Mughees, N., et al. (2026). The
Verifiable Action Card: Trustworthy Human-in-the-Loop Control for Secure
Autonomous Agents. arXiv. \url{https://doi.org/10.48550/arXiv.2609.18411}

\bibitem{loopjacking} Kumar, A. A. (2026). Loopjacking: Hijacking
Human-in-the-Loop Approval. arXiv.
\url{https://doi.org/10.48550/arXiv.2609.21081}

\bibitem{oversight} Holzinger, A., Zatloukal, K., Müller, H. (2024). Is human
oversight to AI systems still possible? New Biotechnology.
\url{https://doi.org/10.1016/j.nbt.2024.12.003}

\end{thebibliography}

\end{document}
